% Lista transcrita a partir de HTML do Projeto Agatha.
% Origem: Razão, Proporção e Porcentagem I
%
% FIGURAS MANTIDAS:
% Q08, Q15, Q16, Q20, Q21, Q22, Q25
%
% Q21:
% As alternativas eram imagens com expressões matemáticas e foram reconstruídas em LaTeX.

---
tipo: Múltipla Escolha
dificuldade: Média
nivel: médio
disciplina: Matemática
assunto: Razões e Proporções
gabarito: A
resposta:
tags: []
fonte: original
concurso: ENEM
banca: INEP
ano: 2025
numero: "01"
---
\question
A cada bimestre, a diretora de uma escola compra uma quantidade de folhas de papel ofício proporcional ao número de alunos matriculados. No bimestre passado, ela comprou 6 000 folhas para serem utilizadas pelos 1 200 alunos matriculados. Neste bimestre, alguns alunos cancelaram suas matrículas e a escola tem, agora, 1 150 alunos. A diretora só pode gastar R$ 220,00 nessa compra, e sabe que o fornecedor da escola vende as folhas de papel ofício em embalagens de 100 unidades a R$ 4,00 a embalagem. Assim, será preciso convencer o fornecedor a dar um desconto à escola, de modo que seja possível comprar a quantidade total de papel ofício necessária para o bimestre.

O desconto necessário no preço final da compra, em porcentagem, pertence ao intervalo
\begin{choices}
\CorrectChoice (5,0 ; 5,5)
\choice (8,0 ; 8,5)
\choice (11,5 ; 12,5)
\choice (19,5 ; 20,5)
\choice (3,5 ; 4,0)
\end{choices}

---
tipo: Múltipla Escolha
dificuldade: Média
nivel: médio
disciplina: Matemática
assunto: Sistemas Lineares
gabarito: B
resposta:
tags: []
fonte: original
concurso: ENEM
banca: INEP
ano: 2025
numero: "02"
---
\question
O metrô de um município oferece dois tipos de tiquetes com colorações diferentes, azul e vermelha, sendo vendidos em cartelas, cada qual com nove tiquetes da mesma cor e mesmo valor unitário. Duas cartelas de tiquetes azuis e uma cartela de tiquetes vermelhos são vendidas por R$ 32,40. Sabe-se que o preço de um tiquete azul menos o preço de um tiquete vermelho é igual ao preço de um tiquete vermelho mais cinco centavos.

Qual o preço, em real, de uma cartela de tiquetes vermelhos?
\begin{choices}
\choice 4,68
\CorrectChoice 6,30
\choice 9,30
\choice 10,50
\choice 10,65
\end{choices}

---
tipo: Múltipla Escolha
dificuldade: Média
nivel: médio
disciplina: Matemática
assunto: Porcentagem
gabarito: D
resposta:
tags: []
fonte: original
concurso: ENEM
banca: INEP
ano: 2025
numero: "03"
---
\question
Um pescador tem um custo fixo diário de R$ 900,00 com combustível, iscas, manutenção de seu barco e outras pequenas despesas. Ele vende cada quilograma de peixe por R$ 5,00. Souza, o melhor obter um lucro mínimo de R$ 800,00 por dia. Sozinho, ele consegue, ao final de um dia de trabalho, pescar 180 kg de peixe, o que é suficiente apenas para cobrir o custo fixo diário. Portanto, precisa contratar ajudantes, pagando para cada um R$ 250,00 por dia de trabalho. Além desse valor, 4% da receita obtida pela venda de peixe é repartida igualmente entre os ajudantes. Considerando o tamanho de seu barco, ele pode contratar até 5 ajudantes. Ele sabe que com um ajudante a pesca diária é de 300 kg e que, a partir do segundo ajudante contratado, aumenta-se em 100 kg a quantidade de peixe pescada por ajudante em um dia de trabalho.

A quantidade mínima de ajudantes que esse pescador precisa contratar para conseguir o lucro diário pretendido é
\begin{choices}
\choice 1
\choice 2
\choice 3
\CorrectChoice 4
\choice 5
\end{choices}

---
tipo: Múltipla Escolha
dificuldade: Média
nivel: médio
disciplina: Matemática
assunto: Porcentagem
gabarito: A
resposta:
tags: []
fonte: original
concurso: ENEM
banca: INEP
ano: 2025
numero: "04"
---
\question
Dirigir após ingerir bebidas alcoólicas é uma atitude extremamente perigosa, uma vez que, a partir da primeira dose, a pessoa já começa a ter perda de sensibilidade de movimentos e de reflexos. Apesar de a eliminação do álcool depender de cada pessoa e de como o organismo consegue metabolizar a substância, ao final da primeira hora após a ingestão, a concentração de álcool (C) no sangue corresponde a aproximadamente 90% da quantidade (q) de álcool ingerida, e a eliminação total dessa concentração pelo fígado ocorre até 12 horas.

\credits{Disponível em: http://g1.globo.com Acesso em: 1 dez. 2018 (adaptado).}

Nessas condições, ao final da primeira hora após a ingestão da quantidade q de álcool, a concentração C dessa substância no sangue é expressa algebricamente por
\begin{choices}
\CorrectChoice C = 0,9q
\choice C = 0,1q
\choice C = 1 - 0,1q
\choice C = 1 - 0,9q
\choice C = q - 10
\end{choices}

---
tipo: Múltipla Escolha
dificuldade: Média
nivel: médio
disciplina: Matemática
assunto: Razões e Proporções
gabarito: E
resposta:
tags: []
fonte: original
concurso: ENEM
banca: INEP
ano: 2025
numero: "05"
---
\question
Um supermercado conta com cinco caixas disponíveis para pagamento. Foram instaladas telas que apresentam o tempo médio gasto por cada caixa para iniciar e finalizar o atendimento de cada cliente, e o número de pessoas presentes na fila de cada caixa em tempo real. Um cliente, na hora de passar sua compra, sabendo que cada um dos cinco caixas iniciará um novo atendimento naquele momento, pretende gastar o menor tempo possível de espera na fila. Ele observa que as telas apresentam as informações a seguir:

\begin{itemize}
\item Caixa I: atendimento 12 minutos, 5 pessoas na fila.
\item Caixa II: atendimento 6 minutos, 9 pessoas na fila.
\item Caixa III: atendimento 5 minutos, 6 pessoas na fila.
\item Caixa IV: atendimento 15 minutos, 2 pessoas na fila.
\item Caixa V: atendimento 9 minutos, 3 pessoas na fila.
\end{itemize}

Para alcançar seu objetivo, o cliente deverá escolher o caixa
\begin{choices}
\choice I
\choice II
\choice III
\choice IV
\CorrectChoice V
\end{choices}

---
tipo: Múltipla Escolha
dificuldade: Média
nivel: médio
disciplina: Matemática
assunto: Razões e Proporções
gabarito: D
resposta:
tags: []
fonte: original
concurso: ENEM
banca: INEP
ano: 2025
numero: "06"
---
\question
Alguns estudos comprovam que os carboidratos fornecem energia ao corpo, preservam as proteínas estruturais dos músculos durante a prática de atividade física e ainda dão força para o cérebro coordenar os movimentos, o que de fato tem impacto positivo no desenvolvimento do praticante. o ideal é consumir 1 grama de carboidrato para cada minuto de caminhada.

\credits{CIRINO. C. Boa pergunta: consumir carboidratos antes dos exercícios melhora o desempenho do atleta? Revista Saúde! É Vital n. 330, nov. 2010 (adaptado).}

Um casal realizará diariamente 30 minutos de caminhada, ingerindo, antes dessa atividade, a quantidade ideal de carboidratos recomendada. Para ter o consumo ideal apenas por meio do consumo de pão de fôrma integral, o casal planeja garantir o suprimento de pães para um período de 30 dias ininterruptos. Sabe-se que cada pacote desse pão vem com 18 fatias, e que cada uma delas tem 15 gramas de carboidratos.

A quantidade mínima de pacotes de pão de fôrma necessários para prover o suprimento a esse casal é
\begin{choices}
\choice 1
\choice 4
\choice 6
\CorrectChoice 7
\choice 8
\end{choices}

---
tipo: Múltipla Escolha
dificuldade: Média
nivel: médio
disciplina: Matemática
assunto: Razões e Proporções
gabarito: B
resposta:
tags: []
fonte: original
concurso: ENEM
banca: INEP
ano: 2025
numero: "07"
---
\question
Entre maratonistas, um parâmetro utilizado é o de economia de corrida (EC). O valor desse parâmetro é calculado pela razão entre o consumo de oxigênio, em mililitro (mL) por minuto (min), e a massa, em quilograma (kg), do atleta correndo a uma velocidade constante.

\credits{Disponível em: www.treinamentoonline.com.br. Acesso em: 23 out. 2019 (adaptado).}

Um maratonista, visando melhorar sua performance, auxiliado por um médico, mensura o seu consumo de oxigênio por minuto à velocidade constante. Com base nesse consumo e na massa do atleta, o médico calcula o EC do atleta.

A unidade de medida da grandeza descrita pelo parâmetro EC é
\begin{choices}
\choice min/mL * kg
\CorrectChoice mL/min * kg
\choice min * mL/kg
\choice min * kg/mL
\choice mL * kg/min
\end{choices}

---
tipo: Múltipla Escolha
dificuldade: Média
nivel: médio
disciplina: Matemática
assunto: Razões e Proporções
gabarito: E
resposta:
tags: []
fonte: original
concurso: ENEM
banca: INEP
ano: 2025
numero: "08"
---
\question
Um controlador de voo dispõe de um instrumento que descreve a altitude de uma aeronave em voo, em função da distância em solo. Essa distância em solo é a medida na horizontal entre o ponto de origem do voo até o ponto que representa a projeção ortogonal da posição da aeronave, em voo, no solo. Essas duas grandezas são dadas numa mesma unidade de medida.

A tela do instrumento representa proporcionalmente as dimensões reais das distâncias associadas ao voo. A figura apresenta a tela do instrumento depois de concluída a viagem de um avião, sendo a medida do lado de cada quadradinho da malha igual a 1 cm.

\includegraphics{https://mpfaraujo.com.br/guardafiguras/uploads/3d19c3aeafb87f0fa960b4ea4eda8cc6d5b89adda60993098508f7a285ff129.webp}

Essa tela apresenta os dados de um voo cuja maior altitude alcançada foi de 5 km.

A escala em que essa tela representa as medidas reais é
\begin{choices}
\choice 1:5
\choice 1:11
\choice 1:55
\choice 1:5 000
\CorrectChoice 1:500 000
\end{choices}

---
tipo: Múltipla Escolha
dificuldade: Média
nivel: médio
disciplina: Matemática
assunto: Razões e Proporções
gabarito: C
resposta:
tags: []
fonte: original
concurso: ENEM
banca: INEP
ano: 2025
numero: "09"
---
\question
O calendário maia apresenta duas contagens simultâneas de anos, o chamado ano Tzolk'in, composto por 260 dias e que determinava o calendário religioso, e o ano Haab, composto por 365 dias e que determinava o calendário agrícola. Um historiador encontrou evidências de que gerações de uma mesma família governaram certa comunidade maia pelo período de 20 ciclos, sendo cada ciclo formado por 52 anos Haab.

\credits{Disponível em: www.suapesquisa.com Acesso em: 20 ago 2014.}

De acordo com as informações fornecidas, durante quantos anos Tzolk'in aquela comunidade maia foi governada por tal família?
\begin{choices}
\choice 741
\choice 1 040
\CorrectChoice 1 460
\choice 2 100
\choice 5 200
\end{choices}

---
tipo: Múltipla Escolha
dificuldade: Média
nivel: médio
disciplina: Matemática
assunto: Razões e Proporções
gabarito: D
resposta:
tags: []
fonte: original
concurso: ENEM
banca: INEP
ano: 2025
numero: "10"
---
\question
Um agricultor é informado sobre um método de proteção para sua lavoura que consiste em inserir larvas específicas, de rápida reprodução. A reprodução dessas larvas faz com que sua população multiplique-se por 10 a cada 3 dias e, para evitar eventuais desequilíbrios, é possível cessar essa reprodução aplicando-se um produto X. O agricultor decide iniciar esse método com 100 larvas e dispõe de 5 litros do produto X, cuja aplicação recomendada é de exatamente 1 litro para cada população de 200 000 larvas. A quantidade total do produto X de que ele dispõe deverá ser aplicada de uma única vez.

Quantos dias após iniciado esse método o agricultor deverá aplicar o produto X?
\begin{choices}
\choice 2
\choice 4
\choice 6
\CorrectChoice 12
\choice 18
\end{choices}

---
tipo: Múltipla Escolha
dificuldade: Média
nivel: médio
disciplina: Matemática
assunto: Razões e Proporções
gabarito: B
resposta:
tags: []
fonte: original
concurso: ENEM
banca: INEP
ano: 2025
numero: "11"
---
\question
Um borrifador de atuação automática libera, a cada acionamento, uma mesma quantidade de inseticida. O recipiente desse produto, quando cheio, contém 360 mL de inseticida, que duram 60 dias se o borrifador permanecer ligado ininterruptamente e for acionado a cada 48 minutos.

A quantidade de inseticida que é liberada a cada acionamento do borrifador, em mililitro, é
\begin{choices}
\choice 0,125
\CorrectChoice 0,200
\choice 4,800
\choice 6,000
\choice 12,000
\end{choices}

---
tipo: Múltipla Escolha
dificuldade: Média
nivel: médio
disciplina: Matemática
assunto: Porcentagem
gabarito: D
resposta:
tags: []
fonte: original
concurso: ENEM PPL
banca: INEP
ano: 2021
numero: "12"
---
\question
Um ciclista faz um treino para uma prova, em um circuito oval, cujo percurso é de 800 m. Nesse treino, realiza 20 voltas. Ele divide seu treino em 3 etapas. Na primeira etapa, inicializa seu cronômetro e realiza as cinco primeiras voltas com velocidade média de 4 m/s. Na segunda etapa, faz mais cinco voltas, mas com velocidade média 25\% maior que a da etapa anterior. Na última etapa, finaliza o treino mantendo a velocidade média da primeira etapa.

Ao final do treino, o cronômetro estará marcando, em
\begin{choices}
\choice 2 600
\choice 2 800
\choice 3 000
\CorrectChoice 3 800
\choice 4 000
\end{choices}

---
tipo: Múltipla Escolha
dificuldade: Média
nivel: médio
disciplina: Matemática
assunto: Porcentagem
gabarito: C
resposta:
tags: []
fonte: original
concurso: ENEM
banca: INEP
ano: 2025
numero: "13"
---
\question
Um médico faz o acompanhamento clínico de um grupo de pessoas que realizam atividades físicas diariamente. Ele observou que a perda média de massa dessas pessoas para cada hora de atividade física era de 1,5 kg. Sabendo que a massa de 1 L de água é de 1 kg, ele recomendou que ingerissem, ao longo das 3 horas seguintes ao final da atividade, uma quantidade total de água correspondente a 40\% a mais do que a massa perdida na atividade física, para evitar desidratação.

Seguindo a recomendação médica, uma dessas pessoas ingeriu, certo dia, um total de 1,7 L de água após terminar seus exercícios físicos.

Para que a recomendação médica tenha efetivamente sido respeitada, a atividade física dessa pessoa, nesse dia, durou
\begin{choices}
\choice 30 minutos ou menos
\choice mais de 35 e menos de 45 minutos
\CorrectChoice mais de 45 e menos de 55 minutos
\choice mais de 60 e menos de 70 minutos
\choice 70 minutos ou mais
\end{choices}

---
tipo: Múltipla Escolha
dificuldade: Média
nivel: médio
disciplina: Matemática
assunto: Razões e Proporções
gabarito: A
resposta:
tags: []
fonte: original
concurso: ENEM
banca: INEP
ano: 2025
numero: "14"
---
\question
Definem-se o dia e o ano de um planeta de um sistema solar como sendo, respectivamente, o tempo que o planeta leva para dar 1 volta completa em tomo de seu próprio eixo de rotação e o tempo para dar 1 volta completa em tomo de seu Sol.

Suponha que exista um planeta Z, em algum sistema solar, onde um dia corresponda a 73 dias terrestres e que 2 de seus anos correspondam a 1 ano terrestre. Considere que 1 ano terrestre tem 365 de seus dias.

No planeta Z, seu ano corresponderia a quantos de seus dias?
\begin{choices}
\CorrectChoice 2,5
\choice 10,0
\choice 730,0
\choice 13 322,5
\choice 53 290,0
\end{choices}

---
tipo: Múltipla Escolha
dificuldade: Média
nivel: médio
disciplina: Matemática
assunto: Razões e Proporções
gabarito: A
resposta:
tags: []
fonte: original
concurso: ENEM
banca: INEP
ano: 2025
numero: "15"
---
\question
Uma pessoa precisa contratar um operário para fazer um serviço em sua casa. Para isso, ela postou um anúncio em uma rede social.

Cinco pessoas responderam informando preços por hora trabalhada, gasto diário com transporte e tempo necessário para conclusão do serviço, conforme valores apresentados no quadro.

\includegraphics{https://mpfaraujo.com.br/guardafiguras/uploads/54861475cf4af1298943bf36c1777518a1b1b830867662ff11def54081211f58.webp}

Se a pessoa pretende gastar o mínimo possível com essa contratação, irá contratar o operário
\begin{choices}
\CorrectChoice I
\choice II
\choice III
\choice IV
\choice V
\end{choices}

---
tipo: Múltipla Escolha
dificuldade: Média
nivel: médio
disciplina: Matemática
assunto: Razões e Proporções
gabarito: D
resposta:
tags: []
fonte: original
concurso: ENEM
banca: INEP
ano: 2025
numero: "16"
---
\question
Uma instituição de ensino superior ofereceu vagas em um processo seletivo de acesso a Seus cursos. Finalizadas as inscrições, foi divulgada a relação do número de candidatos por vaga em cada um dos cursos oferecidos. Esses dados são apresentados no quadro.

\includegraphics{https://mpfaraujo.com.br/guardafiguras/uploads/42c0c0f65077759e1e18280d65160e22ae65da31afdc6a3066218775ecc8a63e.webp}

Qual foi o número total de candidatos inscritos nesse processo seletivo?
\begin{choices}
\choice 200
\choice 400
\choice 1 200
\CorrectChoice 1 235
\choice 7 200
\end{choices}

---
tipo: Múltipla Escolha
dificuldade: Média
nivel: médio
disciplina: Matemática
assunto: Razões e Proporções
gabarito: E
resposta:
tags: []
fonte: original
concurso: ENEM
banca: INEP
ano: 2025
numero: "17"
---
\question
O pacote básico de um jogo para smartphone, que é vendido a R$ 50,00, contém 2 000 gemas e 100 000 moedas de ouro, que são itens utilizáveis nesse jogo.

A empresa que comercializa esse jogo decidiu criar um pacote especial que será vendido a R$ 100,00 e que se diferenciará do pacote básico por apresentar maiores quantidades de gemas e moedas de ouro. Para estimular as vendas desse novo pacote, a empresa decidiu inserir nele 6 000 gemas a mais, em relação ao que o cliente teria caso optasse por comprar, com a mesma quantia, dois pacotes básicos.

A quantidade de moedas de ouro que a empresa deverá inserir ao pacote especial, para que seja mantida a mesma proporção existente entre as quantidades de gemas e de moedas de ouro contidas no pacote básico, é
\begin{choices}
\choice 50 000
\choice 10 0000
\choice 200 000
\choice 300 000
\CorrectChoice 400 000
\end{choices}

---
tipo: Múltipla Escolha
dificuldade: Média
nivel: médio
disciplina: Matemática
assunto: Razões e Proporções
gabarito: E
resposta:
tags: []
fonte: original
concurso: ENEM PPL
banca: INEP
ano: 2021
numero: "18"
---
\question
Uma loja que vende tintas tem uma máquina que efetua misturas de variadas cores para obter diferentes tonalidades. Um cliente havia comprado 7 litros de tinta de uma tonalidade, proveniente da mistura das cores verde e branco, na proporção de 5 para 2, respectivamente. Tendo sido insuficiente a quantidade de tinta comprada, o cliente retorna à loja para comprar mais 3,5 litros da mesma mistura de tintas, com a mesma tonalidade que havia comprado anteriormente.

A quantidade de tinta verde, em litro, que o funcionário dessa loja deverá empregar na mistura com a tinta branca para conseguir a mesma tonalidade obtida na primeira compra é
\begin{choices}
\choice 14
\choice 1,5
\choice 1,7
\choice 2,3
\CorrectChoice 2,5
\end{choices}

---
tipo: Múltipla Escolha
dificuldade: Média
nivel: médio
disciplina: Matemática
assunto: Razões e Proporções
gabarito: B
resposta:
tags: []
fonte: original
concurso: ENEM PPL
banca: INEP
ano: 2021
numero: "19"
---
\question
Um técnico gráfico constrói uma nova folha a partir das medidas de uma folha AO. As medidas de uma folha AO são 595 mm de largura e 840 mm de comprimento. A nova folha foi construída do seguinte modo: acrescenta uma polegada na medida da largura e 16 polegadas na medida do comprimento. Esse técnico precisa saber a razão entre as medidas da largura e do comprimento, respectivamente, dessa nova folha.

Considere 2,5 cm como valor aproximado para uma polegada.

Qual é a razão entre as medidas da largura e do comprimento da nova folha?
\begin{choices}
\choice 1/16
\CorrectChoice 620/1 240
\choice 596/856
\choice 598/880
\choice 845/4 840
\end{choices}

---
tipo: Múltipla Escolha
dificuldade: Média
nivel: médio
disciplina: Matemática
assunto: Razões e Proporções
gabarito: C
resposta:
tags: []
fonte: original
concurso: ENEM PPL
banca: INEP
ano: 2021
numero: "20"
---
\question
O rendimento de um carro bicombustível (abastecido com álcool ou gasolina), popularmente conhecido como carro flex, quando abastecido com álcool é menor do que quando abastecido com gasolina, conforme o gráfico, que apresenta o rendimento médio dos carros populares.

\includegraphics{https://mpfaraujo.com.br/guardafiguras/uploads/1b0ec29009f989bb6bdb6858de936f06fd1d867671a1823477563d2c0b56a4ac.webp}

Suponha que um cidadão fez uma viagem, cujo percurso foi de 1 009 km, em um carro popular flex, tendo abastecido o carro nos primeiros 559 km com gasolina e, no restante do percurso, com álcool. Considere que no momento do abastecimento não havia mais combustível no tanque.

Qual o valor mais próximo do rendimento médio do carro ao concluir todo o percurso de 1 009 km?
\begin{choices}
\choice 9,90 km/L
\choice 10,43 km/L
\CorrectChoice 10,84 km/L
\choice 11,00 km/L
\choice 12,11 km/L
\end{choices}

---
tipo: Múltipla Escolha
dificuldade: Média
nivel: médio
disciplina: Matemática
assunto: Razões e Proporções
gabarito: B
resposta:
tags: []
fonte: original
concurso: ENEM PPL
banca: INEP
ano: 2021
numero: "21"
---
\question
Os pneus estão entre os itens mais importantes para a segurança de um carro. Segundo revendedores especializados, o desgaste do pneu em um trajeto é diretamente proporcional ao número de voltas que ele efetua em contato com o solo, sem derrapar, durante esse trajeto, sendo que a constante de proporcionalidade \(k\) depende do material empregado na sua fabricação. O proprietário de um carro, cujo diâmetro do pneu mede \(L\) m, conforme indicado na imagem, pretende obter uma expressão que forneça uma estimativa para a medida do desgaste \(D\) desse pneu ao longo de uma viagem de \(x\) km. Para efeito dos cálculos, considerou o diâmetro do pneu como sendo \(L\), independentemente da extensão do trajeto.

\includegraphics{https://mpfaraujo.com.br/guardafiguras/uploads/4298ac27923c32eb2dbb992acafebdbd44c5b970be82d6e0cdb06d98f6b243c5.webp}

O valor de \(D\) é dado pela expressão
\begin{choices}
\choice \(D=\frac{500.k.x}{\pi.L}\)
\CorrectChoice \(D=\frac{1000.k.x}{\pi.L}\)
\choice \(D=\frac{1000.k.x}{L}\)
\choice \(D=\frac{1000.k.x}{\pi.L^2}\)
\choice \(D=\frac{4000.k.x}{\pi.L^2}\)
\end{choices}

---
tipo: Múltipla Escolha
dificuldade: Média
nivel: médio
disciplina: Matemática
assunto: Razões e Proporções
gabarito: B
resposta:
tags: []
fonte: original
concurso: ENEM PPL
banca: INEP
ano: 2021
numero: "22"
---
\question
Na loteria Lotex, cada aposta corresponde à marcação de cinquenta números em um cartão. Caso o apostador marque uma quantidade inferior a cinquenta números, o sistema completará aleatoriamente a sua aposta até integralizar os cinquenta números necessários. Por exemplo, o cartão de aposta retratado representa as escolhas de um jogador antes que o sistema integralize o seu preenchimento.

\includegraphics{https://mpfaraujo.com.br/guardafiguras/uploads/bc05858addfd66db71acc932759d1f889e46e74dfb345f3395ff4c172930475f.webp}

Com relação ao cartão exibido, o jogador reconhece que o número racional que corresponde ao quociente do número de pontos marcados pelo sistema, em seu jogo, pelo número máximo de pontos para validar a aposta é igual a
\begin{choices}
\choice 11/25
\CorrectChoice 14/25
\choice 14/11
\choice 25/14
\choice 25/11
\end{choices}

---
tipo: Múltipla Escolha
dificuldade: Média
nivel: médio
disciplina: Matemática
assunto: Razões e Proporções
gabarito: C
resposta:
tags: []
fonte: original
concurso: ENEM PPL
banca: INEP
ano: 2022
numero: "23"
---
\question
Com o intuito de fazer bombons para vender, uma doceira comprou uma barra de 2 kg de chocolate e 1 L de creme de leite. De acordo com a receita, cada bombom deverá ter exatamente 34 g de chocolate e 12 mL de creme de leite.

Respeitando os critérios estabelecidos, quantos bombons a doceira poderá fazer utilizando o máximo que puder os ingredientes comprados?
\begin{choices}
\choice 5
\choice 8
\CorrectChoice 58
\choice 71
\choice 83
\end{choices}

---
tipo: Múltipla Escolha
dificuldade: Média
nivel: médio
disciplina: Matemática
assunto: Razões e Proporções
gabarito: A
resposta:
tags: []
fonte: original
concurso: ENEM PPL
banca: INEP
ano: 2022
numero: "24"
---
\question
Um jovem, no trajeto que usa para ir para a escola, sempre passa por um grande relógio digital que há no centro da sua cidade e compara a hora nele mostrada com a hora que marca o seu relógio de pulso. Ao longo de 30 dias de observação, constata que o seu relógio atrasa 2 minutos, a cada 15 dias, em relação ao do centro da cidade.

Após 90 dias, sem nenhum dos dois relógios receberem ajustes e mantida a mesma parcela de atraso diário, ao ler as marcações de horário dos dois relógios, verificou que o do centro da cidade marcava exatamente 7 horas.

Qual horário marcava seu relógio de pulso nesse instante?
\begin{choices}
\CorrectChoice 6 h e 48 min
\choice 6 h e 54 min
\choice 6 h e 58 min
\choice 7 h e 06 min
\choice 7 h e 12 min
\end{choices}

---
tipo: Múltipla Escolha
dificuldade: Média
nivel: médio
disciplina: Matemática
assunto: Razões e Proporções
gabarito: E
resposta:
tags: []
fonte: original
concurso: ENEM
banca: INEP
ano: 2025
numero: "25"
---
\question
Um segmento de reta está dividido em duas partes na proporção áurea quando o todo está para uma das partes na mesma razão em que essa parte está para a outra. Essa constante de proporcionalidade é comumente representada pela letra grega \(\varphi\), e seu valor é dado pela solução positiva da equação \(\varphi^2=\varphi+1\).

Assim como a potência \(\varphi^2\), as potências superiores de \(\varphi\) podem ser expressas da forma \(a\varphi+b\), em que \(a\) e \(b\) são inteiros positivos, como apresentado no quadro.

\includegraphics{https://mpfaraujo.com.br/guardafiguras/uploads/c7a867afb830daaf6469540a718d024603e452bc580a89f69fc5421798a5e10d.webp}

A potência \(\varphi^7\), escrita na forma \(a\varphi+b\) (\(a\) e \(b\) são inteiros positivos), é
\begin{choices}
\choice \(5\varphi+3\)
\choice \(7\varphi+2\)
\choice \(9\varphi+6\)
\choice \(11\varphi+7\)
\CorrectChoice \(13\varphi+8\)
\end{choices}

---
tipo: Múltipla Escolha
dificuldade: Média
nivel: médio
disciplina: Matemática
assunto: Porcentagem
gabarito: C
resposta:
tags: []
fonte: original
concurso: ENEM
banca: INEP
ano: 2025
numero: "26"
---
\question
Uma equipe de marketing digital foi contratada para aumentar as vendas de um produto ofertado em um site de comércio eletrônico. Para isso, elaborou um anúncio que, quando o cliente clica sobre ele, é direcionado para a página de vendas do produto. Esse anúncio foi divulgado em duas redes sociais, A e B, e foram obtidos os seguintes resultados:

\begin{itemize}
\item rede social A: o anúncio foi visualizado por 3 000 pessoas; 10\% delas clicaram sobre o anúncio e foram redirecionadas para o site; 3\% das que clicaram sobre o anúncio compraram o produto. O investimento feito para a publicação do anúncio nessa rede foi de R$ 100,00;
\item rede social B: o anúncio foi visualizado por 1 000 pessoas; 30\% delas clicaram sobre o anúncio e foram redirecionadas para o site; 2\% das que clicaram sobre o anúncio compraram o produto. O investimento feito para a publicação do anúncio nessa rede foi de R$ 200,00.
\end{itemize}

Por experiência, o pessoal da equipe de marketing considera que a quantidade de novas pessoas que verão o anúncio é diretamente proporcional ao investimento realizado, e que a quantidade de pessoas que comprarão o produto também se manterá proporcional à quantidade de pessoas que clicarão sobre o anúncio.

O responsável pelo produto decidiu, então, investir mais R$ 300,00 em cada uma das duas redes sociais para a divulgação desse anúncio e obteve, de fato, o aumento proporcional esperado na quantidade de clientes que compraram esse produto. Para classificar o aumento obtido na quantidade (Q) de compradores desse produto, em consequência dessa segunda divulgação, em relação aos resultados observados na primeira divulgação, o responsável pelo produto adotou o seguinte critério:

\begin{itemize}
\item Q ≤ 60\%: não satisfatório;
\item 60\% < Q ≤ 100\%: regular;
\item 100\% < Q ≤ 150\%: bom;
\item 150\% < Q ≤ 190\%: muito bom;
\item 190\% < Q ≤ 200\%: excelente.
\end{itemize}

O aumento na quantidade de compradores, em consequência dessa segunda divulgação, em relação ao que foi registrado com a primeira divulgação, foi classificado como
\begin{choices}
\choice não satisfatório
\choice regular
\CorrectChoice bom
\choice muito bom
\choice excelente
\end{choices}

---
tipo: Múltipla Escolha
dificuldade: Média
nivel: médio
disciplina: Matemática
assunto: Porcentagem
gabarito: C
resposta:
tags: []
fonte: original
concurso: ENEM
banca: INEP
ano: 2025
numero: "27"
---
\question
Um médico faz o acompanhamento clínico de um grupo de pessoas que realizam atividades físicas diariamente. Ele observou que a perda média de massa dessas pessoas para cada hora de atividade física era de 1,5 kg. Sabendo que a massa de 1 L de água é de 1 kg, ele recomendou que ingerissem, ao longo das 3 horas seguintes ao final da atividade, uma quantidade total de água correspondente a 40\% a mais do que a massa perdida na atividade física, para evitar desidratação.

Seguindo a recomendação médica, uma dessas pessoas ingeriu, certo dia, um total de 1,7 L de água após terminar seus exercícios físicos.

Para que a recomendação médica tenha efetivamente sido respeitada, a atividade física dessa pessoa, nesse dia, durou
\begin{choices}
\choice 30 minutos ou menos
\choice mais de 35 e menos de 45 minutos
\CorrectChoice mais de 45 e menos de 55 minutos
\choice mais de 60 e menos de 70 minutos
\choice 70 minutos ou mais
\end{choices}

---
tipo: Múltipla Escolha
dificuldade: Média
nivel: médio
disciplina: Matemática
assunto: Razões e Proporções
gabarito: D
resposta:
tags: []
fonte: original
concurso: ENEM
banca: INEP
ano: 2025
numero: "28"
---
\question
Um automóvel apresenta um desempenho médio de 16 km/L. Um engenheiro desenvolveu um novo motor a combustão que economiza, em relação ao consumo do motor anterior, 0,1 L de combustível a cada 20 km percorridos.

O valor do desempenho médio do automóvel com o novo motor, em quilômetro por litro, expresso com uma casa decimal, é
\begin{choices}
\choice 15,9
\choice 16,1
\choice 16,4
\CorrectChoice 17,4
\choice 18,0
\end{choices}
